\section{Singular covers and a linear bound for monic Riccati families}
\label{app:singular-cover-riccati}

This section closes a broad first-order template: bounded challenge-degree
identities with constant nonzero quadratic value coefficient have only
linearly many bad nearby labels above agreement $5D/3$, even when the
singular coordinates move with the challenge. The conclusion concerns
\emph{actual solutions of one identity}. Applying it to an entire received
line requires proving that every nearby candidate under consideration
satisfies that identity.

Throughout, $D\ge1$, the characteristic is zero or greater than $D$,
and the domain has $n$ distinct coordinates. Fix a received line $f+zg$
and an integer $D<A\le n$. A solution pair $(z,P)$, $\deg P\le D$, is
\emph{bad} if it has at least $A$ agreements and no degree-at-most-$D$
polynomial agrees with $g$ on its full agreement support. This is precisely
the obstruction to full-support common witnesses.

\begin{theorem}[A fixed cover for singular agreements]
\label{thm:fixed-singular-cover}
Let $Q(X,z,u,v)\ne0$ have total jet degree at most $B\ge1$ and challenge
degree at most $H\ge0$. Suppose a fixed set $S$ of $s$ domain coordinates
and a set of at most $E$ exceptional labels have the following property:
outside those labels, every bad solution of $Q(X,z,P,P')=0$ has at most
$e$ singular agreements outside $S$, where singular means
$Q_v(x,z,P(x),P'(x))=0$. Put
\[
 \kappa=\min\left\{s,\left\lfloor\frac{2s+D}{3}\right\rfloor\right\},
 \qquad L=A-e-\kappa,
\]
and assume $L\ge1$. Define
\[
 \tau=\max\{0,2D-3\},\quad b_0=1+\tau(B-1),\quad
 T_0=b_0+\tau H,\quad K_0=H+BT_0,\quad q_0=B+H,
\]
\[
 C=q_0K_0+q_0T_0\frac{n-D}{A-D}+q_0n.
\]
The number of bad labels is at most
\begin{equation}\label{eq:singular-cover-bound}
 E+\frac nL C+n.
\end{equation}
The last $n$ may be omitted if $\kappa=s$. In particular this is linear
in $n$ for fixed $B,H$, $A-D\ge\eta n$, $L\ge\lambda n$, and $E=O(n)$.
The singular-cover assumption is uniform over the actual candidates;
the separant need not be independent of $P$.
\end{theorem}
\begin{proof}
Theorem~\ref{thm:nonsingular-agreement-mca} bounds labels admitting a bad
witness with at least $L$ nonsingular agreements by $nC/L$.
Outside those and the exceptional labels, select one bad witness per
remaining label. Each has at least
$A-(L-1)-e=\kappa+1$ agreements within $S$. If $\kappa=s$, none exist.
Otherwise, writing $b=\kappa+1\le s$, we have $3b-2s>D$.
Two selected candidates at the same label would agree on at least
$2b-s\ge3b-2s>D$ points and hence coincide. At three distinct labels,
their $S$-agreement sets have an intersection of size at least
$3b-2s>D$. Their polynomials consequently satisfy the same affine
relation as their labels identically. If there are at least three
labels, all selected candidates lie on the affine codeword pencil
$P_z=F_0+zG_0$ determined by the first two; otherwise there are at most
$2\le n$ labels.

For one such pencil, a coordinate at which $F_0=f$ and $G_0=g$ agrees
persistently. Every other coordinate agrees at at most one label.
A bad pair must have at least one such extra agreement, because otherwise
$F_0,G_0$ are full-support common witnesses. There are at most $n$ labels
with an extra agreement. This proves the bound.
\end{proof}

For example, a separant $R_0(X)R_1(X,z)$ with
$\deg_XR_1\le e$ and $\deg_zR_1\le h$ satisfies the hypothesis with
$S$ the domain roots of $R_0$ and $E\le h$. Use one nonzero coefficient
of $R_1$ to exclude labels at which it vanishes identically. A different
incidence argument below allows \emph{all} roots to move.

\begin{theorem}[Bounded-challenge monic Riccati families]
\label{thm:monic-riccati-linear}
Let $T(X,z)\ne0$, and let $T,R,S$ have challenge degree at most a fixed
integer $h\ge0$, with unrestricted $X$-degrees. Consider
\begin{equation}\label{eq:monic-riccati-family}
 T(X,z)P'-P^2+R(X,z)P-S(X,z)=0,\qquad \deg P\le D.
\end{equation}
If $A-D\ge\eta n$ and $A-5D/3\ge\lambda n$ for fixed positive
$\eta,\lambda$, then the number of labels possessing a bad $A$-near
solution is $O_{h,\eta,\lambda}(n)$. The same conclusion holds with
any constant nonzero coefficient of $P^2$ in place of $-1$.
\end{theorem}
\begin{proof}
Write $t=\deg_XT$, with leading coefficient $\tau(z)\ne0$, and
$r=\deg_XR$ (allowing $r=-\infty$).

\emph{Low derivative degree: $t\le2D$.}
Let $S_0$ be the domain coordinates where $T(x,z)$ vanishes identically
in $z$, and let $s=|S_0|$. A nonzero coefficient of $T$ as a polynomial
in $z$ shows $s\le t\le2D$. At each coordinate outside $S_0$, the
polynomial $T(x,z)$ has at most $h$ roots. Thus the total number of
singular coordinate--label incidences outside $S_0$ is at most $hn$.
Set $e=\lfloor\lambda n/2\rfloor$. At most
$hn/(e+1)\le2h/\lambda$ labels have more than $e$ such incidences.
Every other label satisfies the uniform cover hypothesis of
Theorem~\ref{thm:fixed-singular-cover}, since the separant is $T$.
Moreover
\[
 L=A-e-\min\{s,\lfloor(2s+D)/3\rfloor\}
 \ge A-e-5D/3\ge\lambda n/2.
\]
Use $B=2,H=h$ in that theorem.

\emph{High linear degree: $r\ge t>2D$.}
A nonzero coefficient of $S$ above degree $r+D$ restricts solutions to
at most $h$ labels. Otherwise equate coefficients of $X^{r+i}$ in
order $i=D,\ldots,0$. This is a triangular system determining every
coefficient of $P$, with pivot $R_r(z)$. The derivative term involves
only higher message coefficients, and $P^2$ is absent since $r>2D$.
Outside at most $h$ roots of the pivot, Cramer's rule gives one possible
family $P=W/\Delta$, where
\[
 \deg_zW,\deg_z\Delta\le K=h(D+1).
\]
An agreement numerator has degree at most $K+1$ in $z$. If at most $D$
coordinates agree persistently, every nearby label supplies at least
$A-D$ nonpersistent incidences, giving at most $n(K+1)/(A-D)$ labels.
If more than $D$ coordinates persist, interpolation forces an affine
codeword pencil, with at most $n$ bad labels. Further residual equations
can only remove candidates.

\emph{High derivative degree without an identically zero pivot.}
Now $t>2D$ and $r\le t-1$. A nonzero coefficient of $S$ above $t+D-1$
again restricts solutions to at most $h$ labels. Otherwise the top
coefficient pivots are
\[
 d_i(z)=i\tau(z)+R_{t-1}(z),\qquad 1\le i\le D.
\]
If each is a nonzero polynomial, Theorem~\ref{thm:triangular-newton}
applies with nonlinear degree $2D<t$, giving the required linear bound.
Its exceptional set already includes roots of the nonconstant pivots.

\emph{An identically resonant pivot.}
It remains to treat $R_{t-1}=-j\tau$ for some $1\le j\le D$.
The characteristic assumption makes this $j$ unique, and every other
pivot is $(i-j)\tau$. Exclude the roots of $\tau$. Descending coefficient
elimination leaves $a_j=d$ and $a_0=c$ free. At the resonant equation,
the compatibility condition is rational in $z$ and independent of
$c,d$. Its numerator has degree at most $K=h(D+1)$; if nonzero, this
already bounds the labels. Otherwise the remaining equations give
\begin{equation}\label{eq:riccati-resonance-reconstruction}
 P=c+\frac{Q_z(X)+dV_z(X)}{\delta(z)},\qquad
 \delta=\tau^D,\qquad \deg_z(\delta,Q,V)\le K.
\end{equation}
Here $Q_z(0)=V_z(0)=0$, $\deg_XQ\le D$, $\deg_XV=j$, and the leading
coefficient of $V$ is $\delta$. These degree bounds follow by descending
induction: each step multiplies by coefficient polynomials of degree
at most $h$ and introduces at most one factor of $\tau$ in the common
denominator. At most $D$ steps are used; the bound $K$ includes the
compatibility numerator as well.

Translate $X$ by a transcendental element $w$, working over $\F(w)$,
and absorb constant terms into $c$. This preserves all degree bounds
and the original agreement incidences. The translated $Q,V$ again have
zero constant terms; their derivatives at zero are the pretranslation
derivatives at $w$. Since $j\ne0$ in the field,
\[
 v(z)=T(w,z)V'_z(0)\ne0.
\]
Put $M=2K+h$. Multiplication of the residual identity by $\delta^2$
gives coefficient equations of parameter degree at most two in $(c,d)$
and $z$-degree at most $M$. Its constant coefficient, after division by
$\delta$, is
\[
 -\delta c^2+R(w,z)\delta c+T(w,z)Q'_z(0)
       +v(z)d-S(w,z)\delta=0.
\]
Outside the roots of $v$, this solves $d=N(z,c)/v(z)$, where
$\deg_cN=2$, its leading coefficient is $\delta$, and
$\deg_zN\le M$. The total number of roots excluded by $\delta$ and $v$
is at most $K+M$.

Substitute $d=N/v$ in every cleared residual and multiply by $v^2$.
The resulting plane equations in $(z,c)$ have bidegree at most $(3M,4)$.
They are not all zero: the $X^{2j}$ residual before substitution is
\[
 -\delta^2d^2+\alpha(z)c+\beta(z)d+\gamma(z).
\]
There is no $cd$ term because $\deg_XV=j<2j$, and $c^2$ occurs only in
the constant coefficient. Its substituted $c^4$ coefficient is therefore
$-\delta^4$. In particular a vertical common component can occur only
at an already excluded root of $\delta$.

For completeness, the plane incidence bound can be applied directly.
Take the common gcd of all the plane residuals. Its curve components
have total bidegree at most $(3M,4)$. After removing that gcd, a generic
linear combination of the residual quotients has no common component
with one fixed nonzero quotient. Intersection in
$\mathbb P^1\times\mathbb P^1$ bounds the isolated points by $24M$.
Relabel the translated domain coordinates and their received values.
An agreement equation, after substitution, is
\[
 v\bigl(\delta c+Q_z(x)-\delta(f(x)+zg(x))\bigr)+V_z(x)N(z,c)=0.
\]
It has bidegree at most $(J,2)$, where $J=M+K+1$. Thus the total number
of nonpersistent incidences per coordinate on the residual curves is
at most $6M+4J$. On curves with at most $D$ persistent coordinates,
each nearby point supplies at least $A-D$ such incidences. More than
$D$ persistent coordinates force an affine codeword graph by interpolation.
Reconstruction is injective: $c$ is the translated constant coefficient
and $d$ is recovered from the leading coefficient of $P-c-Q/\delta$.
Hence there are at most four such graph components. Each contributes at
most $n$ bad labels. Interpolation at the persistent original domain coordinates, before
translation, gives this pencil over the original field. The resonant case is therefore bounded by
\begin{equation}\label{eq:riccati-resonance-bound}
 K+25M+\frac{n(6M+4J)}{A-D}+4n=O_{h,\eta}(n).
\end{equation}
This proves all cases. Replacing $-P^2$ by a constant nonzero multiple
changes the displayed nonzero leading coefficients, but none of the
degree or incidence estimates.
\end{proof}

At quarter rate the first-order threshold has
$A/D\to4(3+\sqrt{133})/31>5/3$. Thus bounded-challenge monic Riccati
identities cannot supply superlinear bad-label examples there at fixed
agreement slack. This does not cover arbitrary implicit first-order
identities, or equations with a nonconstant quadratic value coefficient.
It also does not prove that an interpolating identity of this form
contains every nearby candidate.

\begin{proposition}[A linear lower example with all agreements singular]
\label{prop:all-singular-linear-lower}
Let $\F_q$ have characteristic greater than $D$, with $n$ distinct domain
coordinates, and choose $D<s<n$. Put $A=s+1$ and fix a subset $S_0$ of
$s$ coordinates. If
\begin{equation}\label{eq:all-singular-direction-existence}
 2^nq^{-(s-D)}(1-1/q)^{-(n-s)}<1,
\end{equation}
there is a received line with no degree-at-most-$D$ direction polynomial
agreeing on $A$ coordinates, but with at least $n-s$ distinct nearby
labels whose selected witnesses have entirely singular agreement supports
for one implicit identity of jet and challenge degree two.
\end{proposition}
\begin{proof}
Write $R_0(X)=\prod_{x\in S_0}(X-x)$. At coordinates in $S_0$ set
$f(x)=0,g(x)=x^D$. At each remaining coordinate $x_i$, choose
$u_i\in\F_q^*$ and set
\[
 f(x_i)=-u_ix_i,\qquad g(x_i)=x_i^D+u_i.
\]
At the label $z=x_i$, the candidate $P_z=zX^D$ agrees exactly on
$S_0\cup\{x_i\}$. Its full support cannot have a common direction:
more than $D$ core coordinates force that direction to be $X^D$, which
fails at $x_i$.

Let $U=v-zDX^{D-1}$ and
\[
 Q(X,z,u,v)=R_0(X)(X-z)U+U^2.
\]
Every displayed candidate satisfies $Q=0$. On it, the separant is
$Q_v=R_0(X)(X-z)$, which vanishes at every agreement coordinate and has
at most one zero outside the fixed set $S_0$. The fixed-cover hypothesis
therefore holds with $e=1$ and no exceptional labels for these witnesses.
Indeed every degree-at-most-$D$ solution is $zX^D+c$: writing
$a=P'-zDX^{D-1}$ gives $a(a+R_0(X)(X-z))=0$, and degree comparison
excludes the second factor. We make no assertion that only $c=0$ can
be nearby for an arbitrary choice of the $u_i$.

To exclude ordinary common agreement on any $A$ coordinates, choose the
$u_i$ independently and uniformly in $\F_q^*$. Subtract $X^D$ from a
prospective direction polynomial. The zero polynomial matches exactly
$s<A$ core coordinates. For a nonzero polynomial $H$ of degree at most
$D$, fix an $A$-subset containing $m$ core coordinates. If $m>D$,
matching is impossible. Otherwise at most $q^{D+1-m}$ polynomials vanish
on the selected core, and each matches the $A-m$ independent outside
values with probability at most $(q-1)^{-(A-m)}$. The probability of any
match on this subset is at most
$q^{D-s}(1-1/q)^{-(n-s)}$. Union bounding over at most $2^n$ subsets
and using~\eqref{eq:all-singular-direction-existence} proves existence.
Without an $A$-agreeing direction polynomial, ordinary common agreement
is impossible regardless of $f$.
\end{proof}

For example, $n=100m$, $D=25m-1$, $s=48m-1$ gives quarter rate,
$A/n=0.48$ above its first-order curve, and $52m+1$ displayed bad labels.
Condition~\eqref{eq:all-singular-direction-existence} holds for sufficiently
large prime fields. Here the fixed-cover theorem has positive linear
slack, so its linear order is sharp. This is a linear lower example,
not a quadratic lower bound for general first-order proximity gaps.
